\documentclass[11pt]{article}
\usepackage{amsmath,amssymb,amsthm,hyperref}
\newtheorem{proposition}{Proposition}
\title{Degree Control in Affine Partial Fractions}
\author{}
\date{November 23, 2026}
\begin{document}
\maketitle

\paragraph{Scope.}
This audit concerns exact rational identities, with bounded polynomial
degrees, on an open real domain. It does not give a lower bound for
strict sign representations on a finite Boolean cube.

\section{The degree-sensitive criterion}

Let $B_1,\ldots,B_H$ be nonconstant affine polynomials in $x_1,\ldots,x_n$.
For a polynomial $f$ of degree $q$, let
$f^{\mathrm h}(z,x)=z^q f(x/z)$ denote its degree-$q$ homogenization,
and let $B_i^{\mathrm h}$ denote the degree-one homogenizations.
For $1\le k\le \min(H,q)$ put
\[
 J_k=\left\langle\prod_{i\in I}B_i^{\mathrm h}:|I|=k\right\rangle
 \subseteq\mathbb R[z,x_1,\ldots,x_n].
\]

\begin{proposition}[Bounded coefficients require the homogeneous ideal]
There are coefficient polynomials $g_I$ of degree at most $q-k$ with
\[
 f=\sum_{|I|=k}g_I\prod_{i\in I}B_i
\]
if and only if $f^{\mathrm h}\in J_k$.
\end{proposition}
\begin{proof}
For the forward implication homogenize every summand to degree $q$,
padding each $g_I$ to degree $q-k$ with powers of $z$. This gives
membership in $J_k$. Conversely, every generator of $J_k$ is
homogeneous of degree $k$. Take the degree-$q$ homogeneous part of an
ideal-membership expression for $f^{\mathrm h}$. Its coefficient
polynomials are homogeneous of degree $q-k$. Setting $z=1$ gives the
required expression and preserves the asserted upper degree bound.
\end{proof}

The corresponding identity for $f/\prod_i B_i$ leaves $H-k$
denominator factors in each summand. In particular, $q=H$ and $k=H-1$
leave one affine denominator and affine numerators. Plain membership
in the dehomogenized ideal is weaker: it can use coefficients of
higher degree, whose leading parts cancel. Dropping flats at infinity
is appropriate for unrestricted affine ideal membership, but does not
by itself preserve the coefficient-degree bound in the proposition.

\section{A concrete admissible-denominator counterexample}

\begin{proposition}[Parallel factors]
Put $B_1=1+x+y$, $B_2=2+x+y$, and $f=xy$.
The affine ideal $\langle B_1,B_2\rangle$ is the whole polynomial ring.
Nevertheless there are no affine polynomials $A_1,A_2$ and constant
$c$ for which
\[
 f=cB_1B_2+A_1B_2+A_2B_1
\]
as a polynomial identity. Both denominator forms have positive
constants and strictly positive slopes on the entire two-dimensional
Boolean cube.
\end{proposition}
\begin{proof}
The ideal contains $B_2-B_1=1$. For the nonexistence assertion,
restrict the claimed identity to $(x,y)=(t,-t)$.
Its left side is $-t^2$, whereas its right side is affine in $t$,
since $B_1=1$ and $B_2=2$ on that line. Evaluating at $t=0,1,-1$
gives a contradiction. The coefficient positivity is immediate.
\end{proof}

The two algebraic assertions are kernel checked as
\texttt{HeadComplexity.parallel\_denominators\_no\_affine\_coefficients}
and \texttt{parallel\_denominators\_quadratic\_coefficients} in
\texttt{Results/AffinePartialFractionDegree.lean}.
The latter verifies the unrestricted identity
$xy=(xy)B_2-(xy)B_1$.
The former includes the optional constant bias after denominator
clearing. Equality of rational functions on an open set would imply
the prohibited polynomial identity; equality only at cube vertices
does not have that implication.

Homogenizing gives $B_1^{\mathrm h}=z+x+y$ and
$B_2^{\mathrm h}=2z+x+y$, whose ideal is
$\langle z,x+y\rangle$. At $(z,x,y)=(0,1,-1)$ both generators vanish
but $f^{\mathrm h}=xy=-1$. Thus the homogeneous criterion detects
exactly the missing condition at infinity.

\section{Primary-source audit}

Claire de Korte and Teresa Yu,
\emph{Partial fraction decompositions on hyperplane arrangements},
\url{https://arxiv.org/html/2602.06531v2}, equation (7), requires
numerator degrees at most $\deg(f)-d$. Its Theorem 5.2 treats the
homogeneous case using vanishing on arrangement flats.
Theorem 5.3, as displayed in both the retrieved HTML and PDF text,
states an affine criterion omitting flats supported at infinity.
For the preceding example, take $H=2$ and $d=1$. The only flat with
at least two factors is supported at infinity, so the displayed
affine vanishing condition is vacuous. The asserted bounded-degree
conclusion would contradict the proposition above. Therefore that
affine statement cannot be imported with equation (7)'s degree bound
without a correction or an additional hypothesis. This observation
does not challenge the homogeneous ideal criterion used above.

The PDF screenshot tool failed, but the primary HTML independently
displays both the degree bound and the affine theorem. No visual PDF
inspection is claimed. This source audit supplies no universal head
upper bound. Selecting admissible factors whose homogeneous ideal
contains a polynomial with any prescribed cube signs remains a
construction problem; merely restating that problem is not progress
on the target conjecture.

\end{document}
